% Worksheets with several tree questions: \def\treemultiplequestions{} before importing.
For the following \ifdefined\treemultiplequestions questions\else question\fi, use this dataclass-based tree. A
\lstinline{Tree[T]} has a \lstinline{label} of type \lstinline{T} and a list of
\lstinline{branches}, each of which is also a \lstinline{Tree[T]}.
A tree is never empty.

\begin{lstlisting}
from __future__ import annotations
from dataclasses import dataclass, field

@dataclass
class Tree[T]:
    """A Tree has a label (of type T) and a list of branches,
    which are trees.

    >>> t = Tree(3, [Tree(1, []), Tree(2, [Tree(1, []), Tree(1, [])])])
    >>> t.branches[0]
    Tree(label=1, branches=[])
    >>> len(t.branches)
    2
    >>> t.branches[1].label
    2
    >>> len(t.branches[1].branches)
    2
    >>> expr = Tree(10, [Tree(6, [Tree(2), Tree('*'), Tree(3)]),
    ...                  Tree('+'), Tree(4)])
    >>> print(expr)
    10
      6
        2
        *
        3
      +
      4
    """
    label: T
    branches: list[Tree[T]] = field(default_factory=list)  # branches defaults to []

    def __str__(self):
        return format_tree(self)

def format_tree(t: Tree, indent='') -> str:
    "Format a tree with each branch indented below its label."
    assert isinstance(t, Tree), f'{t!r} is not a Tree'
    assert isinstance(t.branches, list), f'branches of {t!r} is not a list'
    string = indent + str(t.label)
    for b in t.branches:
        string += '\n' + format_tree(b, indent + '  ')
    return string

def is_leaf(t: Tree) -> bool:
    "Return whether a Tree t is a leaf with no branches."
    return not t.branches
\end{lstlisting}
